\documentclass[10pt,reqno]{amsart}

\usepackage[T1]{fontenc}
\usepackage[utf8]{inputenc}
\usepackage{newtxtext,newtxmath}
\usepackage{microtype}
\usepackage{mathtools}
\usepackage{booktabs}
\usepackage{enumitem}
\usepackage[nameinlink,noabbrev]{cleveref}
\usepackage[a4paper,margin=28mm,headheight=15pt]{geometry}

\setlength{\parskip}{2pt}
\setlist[enumerate]{leftmargin=7mm,itemsep=2pt,topsep=3pt}
\emergencystretch=2em

\newtheorem{theorem}{Theorem}[section]
\newtheorem{proposition}[theorem]{Proposition}
\newtheorem{lemma}[theorem]{Lemma}
\newtheorem{corollary}[theorem]{Corollary}
\theoremstyle{definition}
\newtheorem{definition}[theorem]{Definition}
\newtheorem{remark}[theorem]{Remark}

\title[Curvature-operator blow-up in a circle quotient]
{Unbounded Negative Curvature Operator in a Collapsing Circle Quotient of \(S^4\)}

\author{Argus Open Research Project}
\address{Argus Open Research Project}
\date{August 30, 2026}

\subjclass[2020]{Primary 53C20; Secondary 53C21, 53C30}
\keywords{curvature operator, Cheeger deformation, circle action, fixed point, collapsing metric}

\begin{document}

\begin{abstract}
Consider the unit round sphere
\[
 S^4\subset\mathbb C^2\oplus\mathbb R
\]
with the effective isometric circle action of weights \(1\) and \(2\).
Form the quotient metric on
\(S^4\times_{S^1}(\epsilon S^1)\).
At either fixed pole, we compute the complete curvature operator on
\(\Lambda^2T_pS^4\).  With \(t=\epsilon^{-2}\), its smallest eigenvalue is
\[
 1-2t=1-\frac{2}{\epsilon^2}.
\]
It follows that the curvature operator has no lower bound independent of
\(\epsilon\) as the circle factor collapses.  The sectional curvatures at
the fixed point remain nonnegative; the negative direction is represented by
a nondecomposable two-form, illustrating the distinction between sectional
curvature and the curvature operator in dimension four.
\end{abstract}

\maketitle

\section{The quotient metric}

Write
\[
 S^4=\left\{(z_1,z_2,r)\in\mathbb C^2\oplus\mathbb R:
 |z_1|^2+|z_2|^2+r^2=1\right\}
\]
with the unit round metric \(g\).  Let \(S^1\) act by
\[
 \lambda\cdot(z_1,z_2,r)
 =(\lambda z_1,\lambda^2z_2,r).
\]
This action is isometric and effective.  Its fixed set consists of the two
poles
\[
 p_{\pm}=(0,0,\pm1).
\]

Give the auxiliary circle \(S^1=\mathbb R/(2\pi\mathbb Z)\) the metric
\(\epsilon^2d\theta^2\), and let \(S^1\) act diagonally on
\(S^4\times\epsilon S^1\).  Translation on the second factor makes the
diagonal action free.  The quotient is diffeomorphic to \(S^4\) by
\[
 [x,z]\longmapsto z^{-1}\cdot x.
\]

Let \(K\) be the infinitesimal action field on \(S^4\).  For a tangent vector
\(Y\), the horizontal lift is obtained by minimizing
\[
 \lvert Y+sK\rvert_g^2+\epsilon^2s^2
\]
over \(s\in\mathbb R\).  The minimum occurs at
\[
 s=-\frac{\langle Y,K\rangle}{\lvert K\rvert^2+\epsilon^2}.
\]
Consequently the quotient metric is
\[
 g_{\epsilon}(Y,Z)
 =g(Y,Z)
 -\frac{\langle Y,K\rangle\langle Z,K\rangle}
 {\lvert K\rvert^2+\epsilon^2}.
\]
Putting
\[
 t=\epsilon^{-2}
\]
gives the equivalent expression
\begin{equation}\label{eq:metric}
 g_t=g-\frac{t}{1+t\lvert K\rvert^2}
 K^{\flat}\otimes K^{\flat}.
\end{equation}

\section{The fixed-point two-jet}

Work at \(p=p_+\).  Choose the positively oriented round-orthonormal basis
\[
 e_1=\partial_{x_1},\quad e_2=\partial_{y_1},\quad
 e_3=\partial_{x_2},\quad e_4=\partial_{y_2},
\]
where \(z_j=x_j+iy_j\).  The infinitesimal isotropy
\[
 A=(\nabla K)_p
\]
satisfies
\[
 Ae_1=e_2,\qquad Ae_2=-e_1,\qquad
 Ae_3=2e_4,\qquad Ae_4=-2e_3.
\]

Use \(g\)-normal coordinates obtained from \(\exp_p\).  Equivariance of the
exponential map gives \(K(x)=Ax\) in these coordinates.  Since
\[
 g_{ij}=\delta_{ij}+O(\lvert x\rvert^2),\qquad
 K(0)=0,\qquad \lvert K\rvert^2=O(\lvert x\rvert^2),
\]
the quadratic part of the metric perturbation
\(q=g_t-g\) is
\begin{equation}\label{eq:q}
 q_{ij}=-t(Ax)_i(Ax)_j+O(\lvert x\rvert^4).
\end{equation}
The denominator correction in \eqref{eq:metric} starts in degree four and
does not affect the curvature at \(p\).

Adopt the convention
\[
 R(X,Y)Z=\nabla_X\nabla_YZ-\nabla_Y\nabla_XZ-\nabla_{[X,Y]}Z,
\qquad
 R_{ijkl}=\langle R(e_i,e_j)e_l,e_k\rangle.
\]
The unit round sphere then has curvature operator
\(I_{\Lambda^2}\).  Since \(q(0)=dq(0)=0\), the curvature increment at \(p\)
is
\begin{equation}\label{eq:deltaR}
 \Delta R_{ijkl}
 =\frac12\left(
 q_{jk,il}-q_{jl,ik}-q_{ik,jl}+q_{il,jk}
 \right).
\end{equation}

\begin{lemma}\label{lem:entries}
The nonzero independent entries of the curvature increment are
\[
\begin{aligned}
\Delta R_{1212}&=3t,&
\Delta R_{3434}&=12t,\\
\Delta R_{1234}&=4t,&
\Delta R_{1324}&=2t,&
\Delta R_{1423}&=-2t.
\end{aligned}
\]
\end{lemma}

\begin{proof}
Differentiate \eqref{eq:q} twice:
\[
 q_{ij,k\ell}
 =-t\bigl(A_{ik}A_{j\ell}+A_{i\ell}A_{jk}\bigr).
\]
Substitution into \eqref{eq:deltaR} gives the displayed values.  They satisfy
the first Bianchi identity
\[
 \Delta R_{1234}-\Delta R_{1324}+\Delta R_{1423}=0.
\]
\end{proof}

\section{Bianchi polarization}

Let
\[
 \alpha=e^{12}+2e^{34}.
\]
For a decomposable unit two-form \(\eta=X\wedge Y\), the fixed-point
sectional curvature is
\[
 \sec_{g_t}(X,Y)
 =1+3t\langle\alpha,\eta\rangle^2.
\]
The rank-one operator \(3t\,\alpha\otimes\alpha\) has the same quadratic
values on decomposable two-forms but fails the first Bianchi identity.  With
the selected orientation,
\[
 *e^{12}=e^{34},\qquad
 *e^{13}=-e^{24},\qquad
 *e^{14}=e^{23}.
\]
The required four-form correction is \(-2t*\).

\begin{proposition}\label{prop:operator}
At either fixed pole, the curvature operator is
\[
 \mathcal R_t
 =I+3t\,\alpha\otimes\alpha-2t*.
\]
\end{proposition}

\begin{proof}
The right-hand side has the entries in \Cref{lem:entries}.  The Hodge-star
term vanishes on decomposable quadratic evaluations because
\(\langle*\eta,\eta\rangle=0\), while its alternating contribution restores
the Bianchi identity.
\end{proof}

\section{The complete operator and its spectrum}

In the ordered orthonormal basis
\[
 (e^{12},e^{13},e^{14},e^{23},e^{24},e^{34})
\]
of \(\Lambda^2T_pS^4\), \Cref{prop:operator} gives
\begin{equation}\label{eq:matrix}
\mathcal R_t=
\begin{pmatrix}
1+3t&0&0&0&0&4t\\
0&1&0&0&2t&0\\
0&0&1&-2t&0&0\\
0&0&-2t&1&0&0\\
0&2t&0&0&1&0\\
4t&0&0&0&0&1+12t
\end{pmatrix}.
\end{equation}

\begin{theorem}\label{thm:main}
The smallest eigenvalue of the curvature operator at either fixed pole is
\[
 1-2t=1-\frac{2}{\epsilon^2}.
\]
In particular, the curvature operators of
\(S^4\times_{S^1}(\epsilon S^1)\) are not bounded below uniformly as
\(\epsilon\to0\).
\end{theorem}

\begin{proof}
The eigenvalues of the \((e^{12},e^{34})\) block are
\[
 1+\frac{15+\sqrt{145}}2t,\qquad
 1+\frac{15-\sqrt{145}}2t.
\]
The remaining two \(2\)-by-\(2\) blocks have eigenvalues \(1+2t\) and
\(1-2t\), each with multiplicity two.  Since
\(15-\sqrt{145}>0\), the smallest eigenvalue is \(1-2t\).

Equivalently, the unit two-form
\[
 \omega_-=\frac{e^{13}-e^{24}}{\sqrt2}
\]
is an exact eigenvector:
\[
 \mathcal R_t\omega_-=(1-2t)\omega_-.
\]
As \(t=\epsilon^{-2}\), the eigenvalue tends to \(-\infty\).
\end{proof}

\begin{remark}\label{rem:sectional}
There is no contradiction with the nonnegative sectional-curvature values
at the fixed point.  The two-form \(\omega_-\) is not decomposable.  In
dimension four, a lower bound for the curvature operator is strictly stronger
than a lower bound for sectional curvature.
\end{remark}

\section{Scope}

The example is compact, the circle action is effective and isometric, and
the diagonal quotient is smooth for every \(\epsilon>0\).  It therefore
answers negatively the unrestricted question of whether the curvature
operator remains uniformly bounded below in this collapsing construction.
The argument concerns the curvature operator, not merely one sectional
curvature, and the divergence occurs at a fixed point of codimension four.

\appendix

\section{Checklist for verification}

\begin{enumerate}
  \item Derive the quotient metric \eqref{eq:metric} by horizontal
  minimization and verify \(t=\epsilon^{-2}\).
  \item Check the infinitesimal isotropy matrix with weights \(1\) and \(2\).
  \item Verify that only the quadratic term in \eqref{eq:q} contributes at the
  fixed point.
  \item Recompute the curvature entries from \eqref{eq:deltaR} and check the
  first Bianchi identity.
  \item Verify the polarized formula
  \(I+3t\,\alpha\otimes\alpha-2t*\).
  \item Diagonalize the three \(2\)-by-\(2\) blocks in
  \eqref{eq:matrix}.
  \item Check directly that
  \((e^{13}-e^{24})/\sqrt2\) has eigenvalue \(1-2t\).
\end{enumerate}

\begin{thebibliography}{9}

\bibitem{AIM}
American Institute of Mathematics,
\emph{Flows in Riemannian geometry},
workshop problem list and summary, September 16--20, 2019.

\end{thebibliography}

\end{document}
